\documentclass[11pt,a4paper]{article}

\usepackage[utf8]{inputenc}
\usepackage[T1]{fontenc}
\usepackage{polski}
\usepackage{lmodern}
\usepackage{amsmath, amssymb}
\usepackage[margin=2.3cm]{geometry}
\usepackage{graphicx}
\usepackage{xcolor}
\usepackage{listings}
\usepackage{tcolorbox}
\usepackage{booktabs}
\usepackage{hyperref}
\usepackage{siunitx}
\sisetup{output-decimal-marker={,}}

\hypersetup{colorlinks=true, linkcolor=blue!60!black, urlcolor=blue!60!black}

\lstdefinestyle{matlab}{
  language=Matlab,
  basicstyle=\ttfamily\footnotesize,
  keywordstyle=\color{blue!70!black}\bfseries,
  commentstyle=\color{green!40!black}\itshape,
  stringstyle=\color{red!60!black},
  numbers=left, numberstyle=\tiny\color{gray},
  frame=single, rulecolor=\color{gray!40},
  breaklines=true, showstringspaces=false,
  tabsize=2,
  extendedchars=true, inputencoding=utf8,
  literate={ą}{{\k{a}}}1 {ę}{{\k{e}}}1 {ó}{{\'o}}1
           {ż}{{\.z}}1 {ź}{{\'z}}1 {ć}{{\'c}}1
           {ś}{{\'s}}1 {ń}{{\'n}}1 {ł}{{\l}}1
           {Ą}{{\k{A}}}1 {Ę}{{\k{E}}}1 {Ó}{{\'O}}1
           {Ż}{{\.Z}}1 {Ź}{{\'Z}}1 {Ć}{{\'C}}1
           {Ś}{{\'S}}1 {Ń}{{\'N}}1 {Ł}{{\L}}1,
}

\newtcolorbox{wzor}[1][]{colback=blue!4, colframe=blue!55!black,
  fonttitle=\bfseries, title=#1}

\title{\textbf{Sterowanie pompą zbiornika za pomocą regulatora LQR}\\
       \large Ciągły model w~Simulinku --- dokumentacja projektowa}
\author{Projekt semestralny 6}
\date{\today}

\begin{document}
\maketitle

\section{Opis obiektu}

Obiektem sterowania jest walcowy zbiornik cieczy o~polu podstawy
$A = \pi r^2$, zasilany pompą o~regulowanym przepływie $Q_{in}(t)$.
Odpływ przez zawór u~podstawy zbiornika opisany jest równaniem Torricellego.
Dodatkowy niekontrolowany odpływ (np.\ nieszczelność, pobór wody) modelujemy
jako zakłócenie $d(t)$. Zmienną stanu jest poziom cieczy $h(t)$, wielkością
sterującą $Q_{in}(t)$, a~wyjściem również $h(t)$.

\begin{center}
\begin{tabular}{@{}lcl@{}}
\toprule
Parametr & Symbol & Wartość \\
\midrule
Przyspieszenie ziemskie   & $g$      & $\SI{9.81}{m/s^2}$ \\
Promień zbiornika         & $r$      & $\SI{1.0}{m}$ \\
Pole podstawy             & $A=\pi r^2$ & $\approx \SI{3.1416}{m^2}$ \\
Wysokość zbiornika        & $H$      & $\SI{2.0}{m}$ \\
Promień otworu wypływowego& $r_v$    & $\SI{0.04}{m}$ \\
Pole otworu               & $A_v=\pi r_v^2$ & $\approx \SI{5.03e-3}{m^2}$ \\
Współczynnik wypływu      & $\mu_v$  & $0{,}97$ \\
Poziom zadany             & $h_{ref}$& $\SI{1.0}{m}$ \\
Maksymalny wydatek pompy  & $Q_{pump}$ & $2\,Q_{in,0}$ \\
\bottomrule
\end{tabular}
\end{center}

\section{Model matematyczny}

\subsection{Równanie bilansu masy}

Przy założeniu stałej gęstości cieczy, bilans objętości w~zbiorniku
o~stałym polu przekroju $A$:
\begin{equation}
A\,\dot h(t) \;=\; Q_{in}(t) - Q_{out}(t) - d(t).
\end{equation}

\subsection{Równanie Torricellego --- odpływ}

Prędkość wypływu cieczy przez otwór na wysokości dna wynika z~zasady
zachowania energii (Bernoulli) i~wynosi $v = \sqrt{2gh}$. Z~uwzględnieniem
współczynnika kontrakcji strumienia $\mu_v$:
\begin{equation}
Q_{out}(t) \;=\; \mu_v A_v \sqrt{2\,g\,h(t)}.
\end{equation}

\subsection{Pełny model nieliniowy}

Po podstawieniu:
\begin{wzor}[Model nieliniowy]
\begin{equation}
\dot h(t) \;=\; \frac{1}{A}\Big[\,Q_{in}(t) \;-\; \mu_v A_v \sqrt{2 g\,h(t)} \;-\; d(t)\,\Big].
\label{eq:nonlin}
\end{equation}
\end{wzor}

\section{Punkt pracy}

Stan ustalony ($\dot h=0$, $d=0$, $h=h_{ref}$) wymaga równości dopływu
i~odpływu:
\begin{equation}
Q_{in,0} \;=\; Q_{out,0} \;=\; \mu_v A_v \sqrt{2\,g\,h_{ref}}.
\end{equation}
Dla przyjętych parametrów:
\[
Q_{in,0} = 0{,}97 \cdot \pi\,(0{,}04)^2 \cdot \sqrt{2 \cdot 9{,}81 \cdot 1{,}0}
\approx \SI{0.02160}{m^3/s}.
\]

\section{Linearyzacja --- wyznaczenie macierzy \texorpdfstring{$A,B,C,D$}{A,B,C,D}}

\subsection{Forma przestrzeni stanu}

Ogólną postać ciągłego modelu liniowego zapisujemy jako:
\begin{equation}
\dot{\mathbf x}(t) = \mathbf A\,\mathbf x(t) + \mathbf B\,\mathbf u(t),\qquad
\mathbf y(t) = \mathbf C\,\mathbf x(t) + \mathbf D\,\mathbf u(t),
\end{equation}
gdzie macierze $\mathbf A \in \mathbb R^{n\times n}$, $\mathbf B\in\mathbb R^{n\times m}$,
$\mathbf C\in\mathbb R^{p\times n}$, $\mathbf D\in\mathbb R^{p\times m}$. W~naszym
obiekcie układ jest SISO o~jednym stanie (poziom cieczy), zatem $n=m=p=1$ i~wszystkie
macierze degenerują się do skalarów.

\subsection{Wybór zmiennych stanu, sterowania, wyjścia}

\begin{center}
\begin{tabular}{@{}lll@{}}
\toprule
Oznaczenie & Symbol fizyczny & Odchyłka (stan liniowy) \\
\midrule
Zmienna stanu $x$   & $h$     & $\Delta h = h - h_{ref}$ \\
Sterowanie $u$      & $Q_{in}$& $\Delta Q = Q_{in} - Q_{in,0}$ \\
Wyjście $y$         & $h$     & $\Delta h$ \\
\bottomrule
\end{tabular}
\end{center}

Punkt pracy $(h_{ref},Q_{in,0})$ jest stacjonarny: $f(h_{ref},Q_{in,0}) = 0$,
dzięki czemu w~równaniach liniowych znika wyraz stały.

\subsection{Metoda: rozwinięcie Taylora pierwszego rzędu}

Nieliniową funkcję prawej strony zapisujemy jako:
\begin{equation}
f(h,Q_{in}) \;=\; \frac{1}{A}\Big[Q_{in} - \mu_v A_v\sqrt{2gh}\Big].
\end{equation}
Rozwinięcie wokół $(h_{ref},Q_{in,0})$ daje:
\begin{equation}
f(h,Q_{in}) \;\approx\; \underbrace{f(h_{ref},Q_{in,0})}_{=0}
+ \left.\frac{\partial f}{\partial h}\right|_{0}\!\Delta h
+ \left.\frac{\partial f}{\partial Q_{in}}\right|_{0}\!\Delta Q.
\end{equation}
Wyliczenie pochodnych cząstkowych:
\begin{align}
\left.\frac{\partial f}{\partial h}\right|_{0}
&= \frac{1}{A}\cdot\left(-\mu_v A_v\cdot\frac{d}{dh}\sqrt{2gh}\right)\bigg|_{h_{ref}}
 = \frac{1}{A}\cdot\left(-\mu_v A_v\cdot\frac{2g}{2\sqrt{2gh_{ref}}}\right) \notag \\
&= -\,\frac{\mu_v A_v}{A}\,\sqrt{\frac{g}{2\,h_{ref}}}, \\[2pt]
\left.\frac{\partial f}{\partial Q_{in}}\right|_{0} &= \frac{1}{A}.
\end{align}

\subsection{Jawna postać macierzy \texorpdfstring{$A,B,C,D$}{A,B,C,D}}

\begin{wzor}[Macierze modelu liniowego]
\begin{equation}
\mathbf A = \left.\frac{\partial f}{\partial h}\right|_{0}
= -\frac{\mu_v A_v}{A}\sqrt{\frac{g}{2\,h_{ref}}},\qquad
\mathbf B = \left.\frac{\partial f}{\partial Q_{in}}\right|_{0} = \frac{1}{A},
\end{equation}
\begin{equation}
\mathbf C = \left.\frac{\partial h}{\partial h}\right|_{0} = 1,\qquad
\mathbf D = \left.\frac{\partial h}{\partial Q_{in}}\right|_{0} = 0.
\end{equation}
\end{wzor}

\noindent\textbf{Uzasadnienie $\mathbf C,\mathbf D$:} wyjściem jest mierzony
poziom cieczy $y=h$ (tu: odchyłka $\Delta h$). Zatem $y$ zależy wprost od
stanu ($C=1$), a~sterowanie nie wchodzi bezpośrednio do wyjścia ($D=0$), bo
zmiana $Q_{in}$ wpływa na $h$ dopiero przez całkę.

\subsection{Wartości numeryczne}

Podstawiając: $g=9{,}81$, $A=\pi\approx 3{,}1416$, $A_v=\pi\cdot 0{,}04^{2}
\approx 5{,}027\cdot 10^{-3}$, $\mu_v = 0{,}97$, $h_{ref}=1$:
\begin{align}
\mathbf A &= -\frac{0{,}97\cdot 5{,}027\!\cdot\! 10^{-3}}{3{,}1416}
             \sqrt{\tfrac{9{,}81}{2}}
           \approx -\SI{3.437e-3}{s^{-1}}, \\
\mathbf B &= \tfrac{1}{\pi} \approx 0{,}3183, \qquad
\mathbf C = 1,\qquad
\mathbf D = 0.
\end{align}
Stała czasowa obiektu bez regulatora: $T_{ol} = -1/\mathbf A \approx \SI{291}{s}$.
Biegun obiektu leży w~lewej półpłaszczyźnie, więc obiekt jest stabilny,
lecz bardzo powolny --- stąd potrzeba regulatora.

\subsection{Implementacja w~MATLAB-ie}

\begin{lstlisting}[style=matlab]
A_lin = -mu_v*A_v*sqrt(g/(2*h_ref))/A;   % A
B_lin =  1/A;                            % B
C_lin =  1;                              % C
D_lin =  0;                              % D
sys   =  ss(A_lin,B_lin,C_lin,D_lin);    % obiekt liniowy
\end{lstlisting}

\section{Projekt regulatora LQR}

\subsection{Wskaźnik jakości}

Dla układu ciągłego SISO poszukujemy sterowania minimalizującego całkowy
wskaźnik jakości:
\begin{equation}
J \;=\; \int_{0}^{\infty}\!\!\Big(q\,\Delta h^{2}(t) + r\,\Delta Q^{2}(t)\Big)\,dt,
\qquad q=100,\;\; r=1.
\end{equation}

Waga $q$ karze odchyłki poziomu, $r$ --- wysiłek sterowania. Stosunek
$q/r$ reguluje kompromis szybkość--energia.

\subsection{Algebraiczne równanie Riccatiego --- wyprowadzenie}

\subsubsection*{Zasada dynamicznego programowania}

Dla układu ciągłego LTI i kwadratowego wskaźnika jakości postulujemy
\emph{funkcję wartości} (Bellmana) o~formie kwadratowej:
\begin{equation}
V(\mathbf x) \;=\; \mathbf x^{T} \mathbf P\,\mathbf x, \qquad
\mathbf P = \mathbf P^{T} \succ 0.
\end{equation}
Równanie Hamiltona--Jacobiego--Bellmana (HJB):
\begin{equation}
0 = \min_{\mathbf u}\Big\{\,\mathbf x^{T}\mathbf Q\mathbf x
+ \mathbf u^{T}\mathbf R\mathbf u
+ \nabla V^{T}(\mathbf A\mathbf x + \mathbf B\mathbf u)\Big\}.
\end{equation}
Minimum po $\mathbf u$ (przyrównanie pochodnej do zera):
\begin{equation}
\mathbf u^{*} \;=\; -\mathbf R^{-1}\mathbf B^{T}\mathbf P\,\mathbf x
\;\equiv\; -\mathbf K\,\mathbf x.
\label{eq:optu}
\end{equation}
Podstawienie $\mathbf u^{*}$ do HJB i~skrócenie $\mathbf x^{T}(\cdot)\mathbf x$
prowadzi do \textbf{ciągłego algebraicznego równania Riccatiego (CARE)}:

\begin{wzor}[CARE]
\begin{equation}
\mathbf A^{T}\mathbf P + \mathbf P\,\mathbf A
- \mathbf P\,\mathbf B\,\mathbf R^{-1}\mathbf B^{T}\mathbf P
+ \mathbf Q \;=\; \mathbf 0.
\label{eq:care}
\end{equation}
\end{wzor}

\subsubsection*{Wzmocnienie optymalne}

Z~\eqref{eq:optu} bezpośrednio wynika postać wzmocnienia LQR:
\begin{equation}
\mathbf K_{LQR} \;=\; \mathbf R^{-1}\mathbf B^{T}\mathbf P.
\label{eq:Klqr}
\end{equation}
Macierz $\mathbf P$ wybiera się jako \emph{jedyne dodatnio określone}
rozwiązanie~\eqref{eq:care} --- warunek ten gwarantuje stabilność
zamkniętej pętli $(\mathbf A-\mathbf B\mathbf K)$.

\subsection{Zastosowanie CARE do naszego obiektu}

\subsubsection*{Sprowadzenie do wymiaru skalarnego}

Dla $n=m=1$ macierze redukują się do skalarów:
$\mathbf A\to a$, $\mathbf B\to b$, $\mathbf Q\to q$, $\mathbf R\to r$,
$\mathbf P\to P$. Równanie~\eqref{eq:care} przyjmuje postać:
\begin{equation}
a P + P a - P\,b\,r^{-1}\,b\,P + q \;=\; 0
\;\;\Longrightarrow\;\;
\boxed{\,2aP - \frac{b^{2}}{r}\,P^{2} + q = 0\,}.
\label{eq:careS}
\end{equation}

\subsubsection*{Rozwiązanie równania kwadratowego względem \texorpdfstring{$P$}{P}}

Przepisujemy~\eqref{eq:careS} jako $\tfrac{b^{2}}{r}P^{2} - 2aP - q = 0$.
Dyskryminanta:
\begin{equation}
\Delta = (2a)^{2} + 4\tfrac{b^{2}}{r}\,q = 4\Big(a^{2} + \tfrac{b^{2}q}{r}\Big).
\end{equation}
Dwa pierwiastki:
\begin{equation}
P_{\pm} = \frac{2a \pm 2\sqrt{a^{2} + b^{2}q/r}}{2\,b^{2}/r}
        = \frac{r}{b^{2}}\Big(a \pm \sqrt{a^{2} + b^{2}q/r}\Big).
\end{equation}
Ponieważ $\sqrt{a^{2}+b^{2}q/r} > |a|$, tylko pierwiastek ze znakiem
"+" daje $P>0$. Dlatego:
\begin{equation}
\boxed{\,P = \frac{r}{b^{2}}\Big(a + \sqrt{a^{2} + b^{2}q/r}\Big)\,}.
\end{equation}

\subsubsection*{Obliczenie wzmocnienia $K_{LQR}$}

Ze wzoru~\eqref{eq:Klqr}:
\begin{wzor}[Wzmocnienie regulatora]
\begin{equation}
K_{LQR} \;=\; \frac{bP}{r}
       \;=\; \frac{a + \sqrt{a^{2} + b^{2}q/r}}{b}.
\end{equation}
\end{wzor}

\noindent Podstawienie liczbowe ($a=-3{,}437\!\cdot\!10^{-3}$, $b=0{,}3183$,
$q=100$, $r=1$):
\begin{align}
a^{2} &= 1{,}181\cdot 10^{-5}, \\
\tfrac{b^{2}q}{r} &= 0{,}1013 \cdot 100 = 10{,}13, \\
\sqrt{a^{2} + \tfrac{b^{2}q}{r}} &= \sqrt{10{,}13} \approx 3{,}183, \\
K_{LQR} &= \frac{-0{,}003437 + 3{,}183}{0{,}3183} \approx 9{,}989.
\end{align}

\subsection{Weryfikacja --- rozwiązanie CARE w~MATLAB-ie}

Funkcje \texttt{lqr} oraz \texttt{care} rozwiązują~\eqref{eq:care}
numerycznie algorytmem opartym na Hamiltonianie lub dekompozycji Schura:

\begin{lstlisting}[style=matlab]
[K_lqr, P, poles_cl] = lqr(A_lin, B_lin, Q_lqr, R_lqr);
% lub równoważnie:
% [P, poles_cl, K_lqr] = care(A_lin, B_lin, Q_lqr, R_lqr);
\end{lstlisting}
\noindent Wynikiem jest: $P\approx 31{,}37$, $K_{LQR}\approx 9{,}989$,
$\lambda_{cl}\approx -3{,}183$ --- zgodność z~wyprowadzeniem analitycznym.

\subsection{Biegun zamkniętej pętli}

Dla sterowania $\Delta Q = -K_{LQR}\,\Delta h$ równanie zamkniętego układu:
\begin{equation}
\Delta\dot h \;=\; \underbrace{(a - b K_{LQR})}_{\lambda_{cl}}\,\Delta h.
\end{equation}
Podstawienie:
\begin{equation}
\lambda_{cl} = -0{,}003437 - 0{,}3183\cdot 9{,}989 \approx -\SI{3.18}{s^{-1}},
\qquad T_{cl} \approx \SI{0.314}{s}.
\end{equation}
Układ zamknięty jest $\sim$\,925 razy szybszy od obiektu otwartego
($T_{ol}\approx \SI{291}{s}$).

\section{Prawo sterowania}

Po powrocie do~zmiennych fizycznych:
\begin{wzor}[Sterowanie pompy --- model ciągły]
\begin{equation}
\boxed{\;Q_{in}(t) \;=\; Q_{in,0} \;-\; K_{LQR}\,\big(h(t)-h_{ref}\big),\;}
\end{equation}
\end{wzor}
z~saturacją fizyczną pompy $0 \le Q_{in}(t) \le Q_{pump}$.

\section{Schemat modelu w~Simulinku}

Model jest w~pełni \textbf{ciągły} (solver \texttt{ode45}), nieliniowy
(blok \texttt{Sqrt} dla obiektu) oraz ze sprzężeniem zwrotnym poziomu do
regulatora liniowego.

\begin{center}
\begin{tabular}{@{}ll@{}}
\toprule
Blok Simulinka & Funkcja / wzór \\
\midrule
\texttt{Constant} \textsf{h\_ref}         & wartość zadana poziomu \\
\texttt{Sum} $(+,-)$                       & błąd $e = h_{ref}-h$ \\
\texttt{Gain} \textsf{K\_LQR}              & $-K_{LQR}\cdot\Delta h$ \\
\texttt{Constant} \textsf{Q\_in\_0}        & offset punktu pracy \\
\texttt{Sum} $(+,+)$                       & $Q_{in,0}-K_{LQR}(h-h_{ref})$ \\
\texttt{Saturation} $[0,Q_{pump}]$          & ograniczenie pompy \\
\texttt{Sum} $(+,-,-)$                     & $Q_{in}-Q_{out}-d$ \\
\texttt{Gain} $1/A$                        & dzielenie przez pole podstawy \\
\texttt{Integrator}                        & $h(t)=\int\dot h\,dt$, IC $=h_0$ \\
\texttt{Gain}$\cdot$\texttt{Sqrt}$\cdot$\texttt{Gain} & $Q_{out}=\mu_v A_v\sqrt{2gh}$ \\
\texttt{Constant} \textsf{d}               & zakłócenie (domyślnie 0) \\
\texttt{Scope}                             & $h(t)$, $Q_{in}(t)$ \\
\bottomrule
\end{tabular}
\end{center}

\noindent Struktura przepływu sygnałów:
\begin{center}
\small
\texttt{h\_ref} $\to$ $\sum$ $\xrightarrow{e}$ \texttt{K\_LQR} $\to$
$\sum(+Q_{in,0})$ $\to$ sat. $\xrightarrow{Q_{in}}$
$\sum(-Q_{out},-d)$ $\to$ $1/A$ $\to$ $\int$ $\to$ $h(t)$
$\to$ (sprz.\ zwr.) oraz $\to$ $\mu_v A_v\sqrt{2g\cdot}$ $\to$ $Q_{out}$.
\end{center}

\section{Skrypt generujący model}

Model tworzony jest programowo skryptem \texttt{build\_simulink\_lqr.m},
który oblicza parametry $a$, $b$, $K_{LQR}$, wywołuje \texttt{new\_system},
\texttt{add\_block} oraz \texttt{add\_line}, a~następnie zapisuje plik
\texttt{projekt\_lqr\_ciagly.slx}. Uruchomienie:

\begin{lstlisting}[style=matlab]
run build_simulink_lqr.m
sim('projekt_lqr_ciagly')
\end{lstlisting}

\section{Wyniki projektowe (liczbowo)}

\begin{center}
\begin{tabular}{@{}lll@{}}
\toprule
Wielkość & Wzór & Wartość \\
\midrule
$Q_{in,0}$ & $\mu_v A_v\sqrt{2gh_{ref}}$        & \SI{0.02160}{m^3/s} \\
$a$        & $-\tfrac{\mu_v A_v}{A}\sqrt{g/(2h_{ref})}$ & \SI{-3.437e-3}{s^{-1}} \\
$b$        & $1/A$                              & $0{,}3183$ \\
$P$        & $\tfrac{r}{b^2}(a+\sqrt{a^2+b^2q/r})$ & $31{,}37$ \\
$K_{LQR}$  & $bP/r$                             & $9{,}989$ \\
$\lambda_{cl}$ & $a-bK_{LQR}$                   & \SI{-3.18}{s^{-1}} \\
$T_{cl}$   & $-1/\lambda_{cl}$                  & \SI{0.314}{s} \\
\bottomrule
\end{tabular}
\end{center}

\section{Wnioski}

\begin{itemize}\setlength{\itemsep}{2pt}
\item Obiekt jest obiektem pierwszego rzędu nieliniowym ---
      linearyzacja w~punkcie pracy jest wystarczająca, bo LQR utrzymuje
      $h$ blisko $h_{ref}$, gdzie błąd linearyzacji jest mały.
\item Wagi $q=100$, $r=1$ dają bieguny zamkniętej pętli w~$-3{,}18\,\mathrm{s}^{-1}$,
      co odpowiada czasowi regulacji rzędu sekund.
\item Saturacja $[0,Q_{pump}]$ zabezpiecza przed ujemnym przepływem
      (pompa nie wypompowuje wstecznie) oraz przed przekroczeniem
      wydajności fizycznej pompy.
\item W~wersji ciągłej nie używamy histerezy ani przełączania ON/OFF
      obecnych w~oryginalnym skrypcie \texttt{sem6\_projekt\_v1.m};
      pompa pracuje z~ciągle modulowanym wydatkiem.
\end{itemize}

\end{document}
